\chapter{Device Access, AAA and Secure File Transfer}\label{ch:aaa}
\bp{4.1,4.2}

\begin{outcomes}
By the end of this chapter you will be able to:
\begin{tight}
  \item \textbf{Configure} local usernames, enable secrets and SSH-only access to
        the console and VTY lines.
  \item \textbf{Configure} a router or switch as an AAA client against TACACS+ and
        against RADIUS, with a local fallback that cannot lock you out.
  \item \textbf{Compare} TACACS+ and RADIUS and choose correctly between them.
  \item \textbf{Manage} configuration and image files over SCP and SFTP, and say
        why v2.0 replaced TFTP with them.
\end{tight}
\end{outcomes}

\begin{prereq}
\chref{cli} for line configuration and privilege levels. \chref{dns} for
\cmd{ip domain name}, which SSH key generation requires. \chref{transport} for
the difference between TCP and UDP, which is one of the real distinctions between
TACACS+ and RADIUS.
\end{prereq}

\begin{keyterms}
AAA $\bullet$ authentication $\bullet$ authorization $\bullet$ accounting
$\bullet$ TACACS+ $\bullet$ RADIUS $\bullet$ method list $\bullet$ fallback
$\bullet$ \cmd{enable secret} $\bullet$ SCP $\bullet$ SFTP $\bullet$ privilege
level
\end{keyterms}

\begin{examnote}[Two topics, one chapter, and a v2.0 change worth noticing]
Topic~\tp{4.1} is \emph{configure network devices with local usernames and as an
AAA client (TACACS+ and RADIUS) for management}. Topic~\tp{4.2} is \emph{manage
device configuration and software files using secure file transfer operations
with SFTP/SCP}.

They sit together because they are the same job: getting onto a device safely,
and moving files onto it safely. And \tp{4.2} is a genuine v2.0 change --- v1.1
examined TFTP. TFTP is UDP, has no authentication of any kind, and sends
everything in clear text. It has been removed from the blueprint and replaced by
two protocols that fix all three problems. If a practice question offers TFTP as
the secure answer, the question predates your exam.
\end{examnote}

\section{Local access done properly}

Before any server is involved, get the device itself right. Everything in this
section should be on every device you build, including the ones that later use
AAA --- because AAA servers fail, and this is what you fall back to.

\begin{config}{baseline secure access}
hostname CAMPUS-SW1
!
! A domain name is required before RSA keys can be generated: the key
! is named host.domain.
ip domain name campus.example.edu
!
! scrypt is type 9. Use it where the platform offers it. Never use
! "username ... password", which stores type 0 or type 7.
username netadmin privilege 15 algorithm-type scrypt secret S0me-Long-Passphrase
username helpdesk privilege 1  algorithm-type scrypt secret An0ther-Passphrase
!
enable algorithm-type scrypt secret Enable-Passphrase-Here
!
! 2048 is the practical minimum; below 768 the device will refuse to
! run SSH version 2 at all.
crypto key generate rsa modulus 2048
ip ssh version 2
ip ssh time-out 60
ip ssh authentication-retries 3
!
line console 0
 login local
 exec-timeout 10 0
 logging synchronous
!
line vty 0 15
 login local
 transport input ssh
 exec-timeout 10 0
\end{config}

\begin{alertbox}[\cmd{transport input ssh} is the line that matters]
Generating keys and enabling SSH does not disable Telnet. Until you say
\cmd{transport input ssh}, the device still accepts plain-text logins on port~23,
and an auditor will find it. The default on many platforms is
\cmd{transport input all}.

The reverse mistake is worse: typing \cmd{transport input ssh} \emph{before} the
keys exist locks you out of a remote device completely. Generate keys first,
prove SSH works from a second window, then restrict the transport.
\end{alertbox}

\begin{conceptbox}[Password types, and why \cmd{service password-encryption} is not enough]
\begin{tight}
  \item \textbf{Type 0} --- plain text. Visible in the configuration.
  \item \textbf{Type 7} --- Cisco's reversible obfuscation, applied by
        \cmd{service password-encryption}. It is trivially reversed by tools that
        have existed for decades. It stops shoulder-surfing and nothing else.
  \item \textbf{Type 5} --- salted MD5. A real hash, but MD5 is old.
  \item \textbf{Type 8 (PBKDF2) and type 9 (scrypt)} --- current. Use
        \cmd{algorithm-type sha256} or \cmd{algorithm-type scrypt}.
\end{tight}

So \cmd{service password-encryption} is still worth enabling --- it covers things
that can only be type 7, such as line passwords and TACACS+ keys --- but it is
not a substitute for storing user passwords as type 8 or 9.

\cmd{enable secret} is hashed; \cmd{enable password} is not. If both exist, the
secret wins and the password is dead configuration that leaks a credential.
Delete it.
\end{conceptbox}

\section{AAA: three separate questions}

\begin{definitionbox}[AAA]
\term{Authentication} --- who are you? \term{Authorization} --- what are you
allowed to do? \term{Accounting} --- what did you actually do? Three distinct
questions, answered by method lists that a device consults in order.
\end{definitionbox}

\subsection{Method lists}

A method list is an ordered list of ways to answer one of those questions. The
device tries them left to right and moves on only when a method is
\emph{unreachable}, not when it returns a refusal.

\begin{alertbox}[The distinction that saves you and the one that locks you out]
If the TACACS+ server answers ``no, wrong password'', that is an answer. The
device stops there and rejects you --- it does not fall through to the local
database. Fallback happens only when the server does not respond at all.

This is correct behaviour and it surprises people. A fallback method is
protection against a \emph{dead server}, not against a forgotten password.
\end{alertbox}

\begin{config}{AAA against TACACS+, with a local fallback}
aaa new-model
!
! The server. This is the IOS 15 / IOS XE syntax; older configurations
! used a single "tacacs-server host" line.
tacacs server TAC1
 address ipv4 10.0.0.20
 key SharedSecretHere
!
tacacs server TAC2
 address ipv4 10.0.0.21
 key SharedSecretHere
!
aaa group server tacacs+ TACGROUP
 server name TAC1
 server name TAC2
!
! Authentication: try the group, then the local database if no server
! answers. The trailing "local" is not optional in any network you
! intend to keep working.
aaa authentication login default group TACGROUP local
aaa authentication enable default group TACGROUP enable
!
! Authorization: what privilege level does this user get on login?
aaa authorization exec default group TACGROUP local if-authenticated
!
! Accounting: a record of every session and every command at level 15.
aaa accounting exec default start-stop group TACGROUP
aaa accounting commands 15 default start-stop group TACGROUP
!
! Source the packets from a stable address so the server's client
! entry always matches, whichever interface the traffic leaves by.
ip tacacs source-interface Loopback0
\end{config}

\begin{config}{the same device as a RADIUS client}
radius server RAD1
 address ipv4 10.0.0.30 auth-port 1812 acct-port 1813
 key AnotherSharedSecret
!
aaa group server radius RADGROUP
 server name RAD1
!
aaa authentication login default group RADGROUP local
aaa authorization exec default group RADGROUP local
ip radius source-interface Loopback0
\end{config}

\begin{verify}{CAMPUS-SW1}{show aaa servers | include host|state}
RADIUS: id 1, priority 1, host 10.0.0.30, auth-port 1812, acct-port 1813
     State: current UP, duration 43112s, previous duration 0s
\end{verify}

\section{TACACS+ against RADIUS}

Both are AAA protocols and they are not interchangeable. This table is
examinable more or less as it stands.

\begin{longtable}{@{}L{3.5cm}L{5.6cm}L{5.8cm}@{}}
\toprule
\rowcolor{primary!15}
\textbf{} & \textbf{TACACS+} & \textbf{RADIUS} \\
\midrule
\endhead
Transport and port & TCP 49 & UDP 1812 authentication, 1813 accounting \\
\rowcolor{lightbg}
What is encrypted & The entire packet body & The password field only.
Everything else, usernames included, is in clear \\
AAA separation & Fully separate. Authorization can be asked per command &
Authentication and authorization are combined in one exchange \\
\rowcolor{lightbg}
Origin & Cisco-developed, later published & IETF standard, universally
implemented \\
Normal use & \textbf{Device administration} --- controlling who may log in to
routers and switches and what they may type &
\textbf{Network access} --- 802.1X, wireless clients, VPN users \\
\rowcolor{lightbg}
Per-command control & Yes, and this is its main advantage &
No, beyond a privilege level granted at login \\
\bottomrule
\end{longtable}

\begin{examnote}[The one-sentence version]
\textbf{TACACS+ manages administrators; RADIUS manages users.} If a question is
about who may configure a switch, the answer is TACACS+. If it is about a laptop
authenticating onto a wireless network, the answer is RADIUS.

Two supporting facts worth memorising because they are asked directly: TACACS+ is
TCP~49 and encrypts the whole payload; RADIUS is UDP and encrypts only the
password.
\end{examnote}

\section{Moving files securely}

\subsection{Why not TFTP}

TFTP has no user authentication, no encryption, and runs over UDP with no
integrity protection beyond a checksum. Anyone who can reach the server can read
any file whose name they can guess, and a configuration file in flight is
readable by anyone on the path --- including the shared secrets in it. That is
the case for its removal from the blueprint.

\subsection{SCP --- the device as a server}

SCP runs over SSH, so it inherits SSH's authentication and encryption. The device
can act as an SCP \emph{server}, letting you pull and push files from a
workstation.

\begin{config}{enabling the SCP server}
! SCP will not start without all three of these. The authorization
! line in particular is easy to miss, and its absence produces a
! connection that authenticates and then refuses every transfer.
aaa new-model
aaa authentication login default local
aaa authorization exec default local
!
ip scp server enable
\end{config}

With that in place, from a workstation:

\begin{hostcmd}{Linux or macOS}{scp netadmin@10.0.0.1:running-config ./sw1-backup.cfg}
netadmin@10.0.0.1's password:
running-config                       100%  8241     1.2MB/s   00:00
\end{hostcmd}

\subsection{SFTP --- the device as a client}

IOS is an SFTP \emph{client} but not an SFTP server. So the direction of travel
is the device fetching from, or sending to, a server you run.

\begin{verify}{CAMPUS-SW1}{copy sftp: flash:}
Address or name of remote host []? 10.0.0.40
Source username [netadmin]? images
Source filename []? cat9k_iosxe.17.12.04.SPA.bin
Destination filename [cat9k_iosxe.17.12.04.SPA.bin]?
Accessing sftp://images@10.0.0.40/cat9k_iosxe.17.12.04.SPA.bin...
Loading file ! [OK - 1103237120 bytes]
1103237120 bytes copied in 412.338 secs (2675281 bytes/sec)
\end{verify}

\begin{config}{backing the configuration up on a schedule}
! One-off, non-interactive form. The password appears on the command
! line, so prefer key-based SCP where the platform supports it.
! do copy running-config scp://netadmin@10.0.0.40/CAMPUS-SW1.cfg
!
! Automatic archive: a copy every time someone saves, kept off-box.
archive
 path scp://netadmin@10.0.0.40/archive/$h-
 write-memory
 maximum 14
\end{config}

\begin{conceptbox}[Verify the image before you reload onto it]
A transfer that reports \cmd{[OK]} has moved the bytes it was given. It has not
told you the file is the image Cisco shipped. Before reloading:

\begin{tight}
  \item \cmd{verify /md5 flash:cat9k\_iosxe.17.12.04.SPA.bin} and compare against
        the hash published with the download.
  \item \cmd{show flash:} to confirm the size matches.
  \item Confirm free space \emph{before} the copy, not after --- a partial copy
        that filled flash can leave a device unable to boot.
\end{tight}
\end{conceptbox}

\begin{pitfall}
\begin{tight}
  \item \textbf{\cmd{aaa authentication login default group TACGROUP} with no
        \cmd{local}.} When the server is unreachable, nobody can log in --- not
        even at the console. This is the classic lock-out.
  \item \textbf{Testing AAA by closing your session.} Always open a second
        session and prove it works before closing the first.
  \item \textbf{Mismatched shared key.} Authentication fails with no useful
        message on the client. Check the server's logs, not the switch's.
  \item \textbf{No \cmd{ip tacacs source-interface}.} Packets leave by whichever
        interface routing picks, the server does not recognise the source, and
        the failure is intermittent by design.
  \item \textbf{\cmd{transport input ssh} before the keys exist.} Instant remote
        lock-out.
  \item \textbf{Enabling SCP without \cmd{aaa authorization exec}.} Logs in,
        transfers nothing.
  \item \textbf{Believing \cmd{service password-encryption} protects anything.}
        Type 7 is reversible.
  \item \textbf{Reaching for TFTP out of habit.} It is off the blueprint and it
        was off the good-practice list long before that.
\end{tight}
\end{pitfall}

\begin{breakfix}[Nobody can log in to any switch. The AAA server is up and other services work.]
A campus of forty switches, all authenticating to TACACS+. Overnight, logins
begin failing everywhere. The TACACS+ server is running, responds to ping from a
workstation, and is authenticating VPN users normally. Console access fails as
well as SSH.

\begin{verify}{CAMPUS-SW1}{show run | include aaa authentication}
aaa authentication login default group TACGROUP local
aaa authentication enable default group TACGROUP enable
\end{verify}

\begin{verify}{CAMPUS-SW1}{show aaa servers | include host|state}
TACACS+: id 1, priority 1, host 10.0.0.20, auth-port 49
     State: current UP, duration 61204s
\end{verify}

\textbf{Diagnosis.} The configuration has a correct fallback and the server is
reachable and marked \cmd{UP}. That combination rules out both of the usual
suspects at once --- and it is the clue.

Because the server is reachable, the device sends the request and gets an
\emph{answer}. The answer is a rejection. Fallback to \cmd{local} never fires,
because fallback happens on silence, not on refusal. The fault is therefore on
the server: overnight, the group that grants network-device administration was
changed, and the server is now correctly refusing everyone.

The tell is that console access fails too. A network-side or transport-side fault
would leave the console working, because the console does not depend on the
network. A fault that follows you to the console is in the authentication
decision itself.

\textbf{Immediate recovery.} You need an account the switch can validate without
the server. If a local account exists, use it --- but the method list consults
TACACS+ first and TACACS+ is answering, so it will not be reached. Break the
server's reachability from the switch deliberately, or use the password-recovery
procedure on one switch, restore local access, and work from there.

\textbf{The lasting fix} is on the server, not the switches. Restore the
administration group.

\textbf{And the lesson to write down.} A local fallback protects you from an
unreachable server. It does not protect you from a reachable server that says no.
Sites that care about this keep a break-glass local account and test it on a
schedule --- and keep at least one device out of AAA entirely, reachable
out-of-band, as the way back in.
\end{breakfix}

\begin{lab}{Secure access, AAA, and a backup you can restore}{CML or GNS3, with a TACACS+ or RADIUS server}
\textbf{Platform note.} Packet Tracer models a simplified AAA server and will
carry tasks 1--4. Tasks 5--8 need a real server (\cmd{tac\_plus} or FreeRADIUS)
and a real SSH stack, so use CML or GNS3.

\textbf{Topology.} Two switches and a router, a management VLAN, a Linux host
running the AAA server and an SFTP server, and your own workstation.

\begin{tasks}
  \item Build baseline access on all three devices: hostname, domain name, two
        local users at different privilege levels, enable secret, RSA keys, SSH
        version 2, \cmd{transport input ssh}. Verify by SSH and confirm Telnet is
        refused.
  \item Prove the password types. Run \cmd{show run | include username} and
        identify the type number on each. Add one user with
        \cmd{password} instead of \cmd{secret} and compare, then remove it.
  \item Configure AAA against the server with a local fallback. \textbf{Open a
        second SSH session before you test}, and keep it open throughout.
  \item Verify with \cmd{show aaa servers} and by logging in as a server-defined
        user. Confirm the privilege level you land in.
  \item \textbf{Break 1.} Stop the AAA server. Log in again and confirm the local
        fallback works. Record how long the device waits before falling back.
  \item \textbf{Break 2.} Restart the server and disable your test account on it.
        Attempt to log in. Explain why the local fallback does \emph{not} rescue
        you, and relate this to the Break \& Fix above.
  \item \textbf{Break 3.} Remove \cmd{local} from the method list, stop the
        server, and log out of every session. Recover the device. Write down how
        long it took and what you would have done differently.
  \item Enable the SCP server and copy the running configuration to your
        workstation. Then omit \cmd{aaa authorization exec} and observe the
        different failure.
  \item Configure \cmd{archive} to write to SFTP on every \cmd{write memory}.
        Make a change, save, and confirm the archived file appears.
  \item \textbf{Extension.} Copy an image by SFTP, verify its MD5 against the
        published hash, and deliberately truncate a copy to see what
        \cmd{verify} reports.
\end{tasks}
\end{lab}

\begin{checkpoint}
\begin{tightnum}
  \item In \cmd{aaa authentication login default group TACGROUP local}, when is
        \cmd{local} used and when is it not?
  \item Give the transport, port and encryption scope of TACACS+ and of RADIUS,
        and say which you would choose for controlling switch administrators.
  \item Why must \cmd{ip domain name} be set before generating RSA keys?
  \item What three commands does the IOS SCP server require, and what is the
        symptom when the authorization line is missing?
  \item Why did v2.0 replace TFTP with SFTP and SCP? Give three reasons.
\end{tightnum}
\end{checkpoint}

\begin{chaptersummary}
\begin{tight}
  \item Baseline first: hostname, domain name, local users with
        \cmd{algorithm-type scrypt secret}, \cmd{enable secret}, RSA keys,
        \cmd{ip ssh version 2}, \cmd{login local}, \cmd{transport input ssh}.
  \item Type 7 is obfuscation, not encryption. Type 8 and 9 are the current
        choices. \cmd{enable password} should not exist.
  \item AAA answers three separate questions through ordered method lists. A
        method is skipped only when it does not respond --- a rejection ends the
        list.
  \item Always end an authentication method list with \cmd{local}, and always
        test from a second session.
  \item TACACS+: TCP 49, whole packet encrypted, full AAA separation, per-command
        authorization, device administration. RADIUS: UDP 1812/1813, password
        only encrypted, combined authentication and authorization, network
        access.
  \item SCP puts file transfer inside SSH; IOS can be an SCP server and an SFTP
        client. Both replace TFTP, which authenticates nobody and encrypts
        nothing.
  \item Verify an image's hash before reloading onto it.
\end{tight}
\end{chaptersummary}

\begin{reviewq}
  \item \textbf{(MCQ)} Which protocol encrypts the entire packet body and uses
        TCP port 49?
        \begin{tight}
          \item[(a)] RADIUS \quad (b)~TACACS+ \quad (c)~SNMPv3 \quad (d)~SCP
        \end{tight}
  \item \textbf{(MCQ)} A device is configured
        \cmd{aaa authentication login default group SRV local}. The server is
        reachable and rejects the credentials supplied. What happens?
        \begin{tight}
          \item[(a)] The local database is consulted and the user may still log in
          \item[(b)] The login is refused; the local database is not consulted
          \item[(c)] The device retries the server three times, then falls back
          \item[(d)] The user is granted privilege level 1
        \end{tight}
  \item \textbf{(Short)} Explain the difference between \cmd{enable password} and
        \cmd{enable secret}, and say what to do if a configuration contains both.
  \item \textbf{(Short)} Give three concrete reasons SFTP or SCP should be used in
        place of TFTP for moving a configuration file.
  \item \textbf{(Applied)} Write the configuration that makes a switch an AAA
        client of two TACACS+ servers with a local fallback, sourcing its packets
        from \cmd{Loopback0}, and accounting for every level-15 command. Then
        state the two verification steps you would take before ending your
        session.
  \item \textbf{(Applied)} A colleague has enabled the SCP server. Logins
        succeed, transfers fail immediately. Give the missing command and explain
        why authentication alone is not sufficient.
\end{reviewq}
